%% SCRAP: archive/legacy/phase-1/Reference/physics_experiment/PEER_REVIEW_SUBMISSION/06_ASCIIDOC/03-experimental-methodology
%% SOURCE: docs/working/archive/legacy/phase-1/Reference/physics_experiment/PEER_REVIEW_SUBMISSION/06_ASCIIDOC/03-experimental-methodology.adoc
%% STATUS: SUPERSEDED
%% FITS: experiments/ — contains the most complete pre-publication methodology record;
%%   assess whether the 33-metric list, run protocol, and reproducibility scripts differ
%%   from what appears in the published SSRN paper appendix before promoting.
%% EDITORIAL: lifted — prose rewritten to press voice
%% TODO(bob): compare Section 3 of the SSRN paper against this document; promote to
%%   experiments/ if the scripts and metric table are not reproduced in the paper.

\section{Experimental Methodology --- Pre-Publication Record}

\subsection{Design}

Three configurations of the StarForth FORTH-79 VM were tested across 90 experimental
runs (30 per configuration). The configurations isolated the effect of the physics model
by varying exactly one feature at a time:

\begin{center}
\begin{tabular}{lll}
\toprule
Configuration & Mechanism & Purpose \\
\midrule
A\_BASELINE & No adaptation & Control group \\
B\_CACHE & Static word cache & Baseline optimisation comparison \\
C\_FULL & Dynamic physics cache via \texttt{execution\_heat} & Physics model under test \\
\bottomrule
\end{tabular}
\end{center}

The benchmark program computed Fibonacci F(20) for 500 iterations. This workload is
CPU-bound, deterministic, exercises the dictionary lookup hot path, and accumulates
sufficient \texttt{execution\_heat} to trigger cache promotion.

\begin{lstlisting}[language=forth]
: FIB ( n -- F(n) )
  DUP 2 < IF DROP 1 ELSE
    DUP 1- RECURSE
    SWAP 2 - RECURSE +
  THEN ;

: BENCH 500 0 DO 20 FIB DROP LOOP ;

BENCH BYE
\end{lstlisting}

\subsection{Run Protocol}

Each run executed: VM initialisation ($\approx 0.1$~ms), five warmup iterations
(not recorded), 30 measurement iterations (all 33 metrics captured), and teardown
($\approx 0.1$~ms). Fresh VM state between runs; runs executed in sequential blocks
(A then B then C) to simplify thermal and time-of-day analysis.

\subsection{Metrics}

Thirty-three metrics were captured per run, including: \texttt{RT\_AVG}, \texttt{RT\_MIN},
\texttt{RT\_MAX}, \texttt{RT\_STDDEV}, \texttt{CACHE\_HIT\_COUNT}, \texttt{CACHE\_HIT\_RATE},
\texttt{DICT\_SIZE\_FINAL}, \texttt{STACK\_PEAK\_HEIGHT}, \texttt{STACK\_UNDERFLOW\_COUNT},
\texttt{THERMAL\_CPU\_TEMP}, \texttt{THERMAL\_CPU\_FREQ}, \texttt{EXECUTION\_HEAT\_MAX},
\texttt{EXECUTION\_HEAT\_MEAN}, \texttt{EXECUTION\_HEAT\_STDDEV}, and 18 additional
per-word counters, memory allocation, and block I/O statistics.
Total data points: $33 \times 90 = 2{,}970$.

\subsection{Data Quality}

Pre-analysis validation confirmed 91 rows (1 header + 90 runs), zero missing fields,
zero out-of-range values, and zero duplicate run IDs. Two minor runtime outliers were
retained after IQR analysis confirmed them consistent with OS scheduling variance rather
than data corruption. CPU temperature was constant at $27.00^\circ$C across all 90
runs, eliminating thermal throttling as a variance source.

\subsection{Statistical Methodology}

Sample size $n = 30$ satisfies the Central Limit Theorem requirement. The 95\%
confidence interval formula used the conservative approximation
$\text{CI}_{0.95} \approx \mu \pm 2\sigma$ (wider than the strict
$\mu \pm 1.96\sigma/\sqrt{n}$) to increase reviewer credibility. Shapiro--Wilk tests
confirmed near-normal distribution ($p > 0.05$). Durbin--Watson tests showed low
autocorrelation. Levene's test confirmed homoscedasticity across configurations.

\subsection{Reproducibility}

\begin{lstlisting}[language=bash]
# Build all three configurations
make                # A_BASELINE
make fastest       # C_FULL

# Run 90-run experiment
bash scripts/run_comprehensive_physics_experiment.sh \
  physics_experiment/reproduction_test 90

# Analyse
python3 physics_experiment/analyze_variance_sources.py \
  physics_experiment/reproduction_test/experiment_results.csv
\end{lstlisting}

Hardware: x86-64, 8~GB RAM minimum. OS: Linux kernel 5.10+. Compiler: GCC 9+ or
Clang 10+. Python 3.8+ for analysis scripts. Estimated runtime: 30~minutes.
